\documentclass{article}

\usepackage{amsmath}
\usepackage{amsfonts}
\usepackage{amsthm}
\usepackage{amssymb}
\usepackage[margin=2cm]{geometry}
\usepackage{color}
\usepackage{graphicx}

\begin{document}

\title{{\Large IGL Brochure - Exploring Entropy and Complexity of Curves on a Surface}}
\author{Spencer Dowdall (Faculty mentor),
	Joe Nance,
	Victor Yang,
	Xiaolong Han,
	Yohan Kang}
\date{November 14, 2014}
\maketitle

Each closed curve $\gamma$ with a base point $p$ on a surface $S$ determines a natural homeomorphism $\mathcal{P}(\gamma)$ of the surface $S\setminus\{p\}$ via \textbf{point-pushing}:
\begin{figure}[h]
\begin{center}
\includegraphics[width=8cm]{brochure--placeholder_picture}
\caption{{\bf Point-push} around $\gamma$ by dragging basepoint $p$ once around $\gamma$.}
\end{center}
\end{figure}

The goal of our project is to study point-pushing homeomorphisms on the genus-$g$ surface, and to relate their topological entropy to their combinatorial complexity. To this end, we developed a computer program that 
randomly generates closed curves $\gamma$ on the surface, and then calculates the self-intersection number $\iota(\gamma)$ of the curve (measuring the combinatorial complexity) and the dilatation number (measuring the topological entropy) of the resulting point-pushing homeomorphism $\mathcal{P}(\gamma)$. We then statistically analyzed the results by calculating the distributions of the dilatation and self-intersection number data and looking for correlations between these numbers.
\end{document}